\documentclass[12pt,a4paper]{article}

% ---------- Mise en page ----------
\usepackage[a4paper,top=2.7cm,bottom=2.5cm,left=2.3cm,right=2.3cm,headheight=32pt]{geometry}
\usepackage[utf8]{inputenc}
\usepackage[T1]{fontenc}
\usepackage{lmodern}
\usepackage[french]{babel}
\usepackage{amsmath,amssymb,mathtools}
\usepackage{graphicx}
\usepackage[dvipsnames]{xcolor}
\usepackage{titlesec}
\usepackage{enumitem}
\usepackage{fancyhdr}
\usepackage{tcolorbox}
\tcbuselibrary{skins,breakable}
\usepackage{array}
\usepackage{booktabs}
\usepackage{tabularx}
\usepackage{tikz}
\usetikzlibrary{arrows.meta}
\usepackage[hidelinks]{hyperref}

% ---------- Couleurs Mundiapolis ----------
\definecolor{mundiblue}{HTML}{0070C0}
\definecolor{mundidark}{HTML}{123B5E}
\definecolor{mundigray}{HTML}{5A5A5A}
\definecolor{munditeal}{HTML}{1B6B6B}
\definecolor{mundired}{HTML}{B03A2E}

\titleformat{\section}
  {\normalfont\Large\bfseries\color{mundidark}}
  {\thesection}{0.6em}{}
\titleformat{\subsection}
  {\normalfont\large\bfseries\color{mundiblue}}
  {\thesubsection}{0.6em}{}
\titleformat{\subsubsection}
  {\normalfont\normalsize\bfseries\color{mundigray}}
  {\thesubsubsection}{0.6em}{}

\pagestyle{fancy}
\fancyhf{}
\renewcommand{\headrulewidth}{0.6pt}
\renewcommand{\footrulewidth}{0.4pt}
\fancyhead[L]{\small\color{mundigray} Mathématiques appliquées à la gestion -- S1 MGE}
\fancyhead[R]{\small\color{mundigray} Pr Rajae Malek}
\fancyfoot[L]{\small\color{mundigray} Année universitaire 2026--2027}
\fancyfoot[R]{\small\color{mundigray} \thepage}

\newtcolorbox{definitionbox}[1][]{
  breakable, enhanced, colback=blue!4, colframe=mundiblue,
  boxrule=0.6pt, arc=1mm, left=6pt, right=6pt, top=4pt, bottom=4pt,
  fonttitle=\bfseries, coltitle=white, colbacktitle=mundiblue,
  title={Définition #1}
}
\newtcolorbox{propbox}[1][]{
  breakable, enhanced, colback=green!5, colframe=OliveGreen!70!black,
  boxrule=0.6pt, arc=1mm, left=6pt, right=6pt, top=4pt, bottom=4pt,
  fonttitle=\bfseries, coltitle=white, colbacktitle=OliveGreen!70!black,
  title={Propriété #1}
}
\newtcolorbox{exemplebox}[1][]{
  breakable, enhanced, colback=yellow!8, colframe=orange!80!black,
  boxrule=0.6pt, arc=1mm, left=6pt, right=6pt, top=4pt, bottom=4pt,
  fonttitle=\bfseries, coltitle=black, colbacktitle=orange!25,
  title={Exemple appliqué à la gestion -- #1}
}
\newtcolorbox{gestionbox}[1][Application en gestion]{
  breakable, enhanced, colback=orange!6, colframe=orange!80!black,
  boxrule=0.6pt, arc=1mm, left=6pt, right=6pt, top=4pt, bottom=4pt,
  fonttitle=\bfseries, coltitle=black, colbacktitle=orange!30,
  title=#1
}
\newtcolorbox{exobox}[1][]{
  breakable, enhanced, colback=gray!5, colframe=mundidark,
  boxrule=0.6pt, arc=1mm, left=6pt, right=6pt, top=4pt, bottom=4pt,
  fonttitle=\bfseries, coltitle=white, colbacktitle=mundidark,
  title={Énoncé -- #1}
}
\newtcolorbox{corrigebox}[1][]{
  breakable, enhanced, colback=gray!3, colframe=gray!60!black,
  boxrule=0.6pt, arc=1mm, left=6pt, right=6pt, top=4pt, bottom=4pt,
  fonttitle=\bfseries, coltitle=white, colbacktitle=gray!60!black,
  title={Corrigé commenté -- #1}
}
\newtcolorbox{methodebox}[1][]{
  breakable, enhanced, colback=teal!4, colframe=munditeal,
  boxrule=0.6pt, arc=1mm, left=6pt, right=6pt, top=4pt, bottom=4pt,
  fonttitle=\bfseries, coltitle=white, colbacktitle=munditeal,
  title={Méthode -- #1}
}
\newtcolorbox{attentionbox}[1][Piège classique]{
  breakable, enhanced, colback=red!4, colframe=mundired,
  boxrule=0.6pt, arc=1mm, left=6pt, right=6pt, top=4pt, bottom=4pt,
  fonttitle=\bfseries, coltitle=white, colbacktitle=mundired,
  title=#1
}
\newtcolorbox{profbox}[1][Fiche professeur -- quoi faire exactement]{
  breakable, enhanced, colback=mundidark!4, colframe=mundidark,
  boxrule=0.6pt, arc=1mm, left=6pt, right=6pt, top=4pt, bottom=4pt,
  fonttitle=\bfseries, coltitle=white, colbacktitle=mundidark,
  title=#1
}

\setlength{\parindent}{0pt}
\setlength{\parskip}{4pt}
\renewcommand{\arraystretch}{1.25}

\begin{document}

% =======================================================
\begin{titlepage}
\hypersetup{pageanchor=false}
\centering
\vspace*{0.3cm}
\IfFileExists{logo.png}{%
  \includegraphics[width=5.2cm]{logo.png}\\[1.2cm]
}{%
  {\Huge\bfseries\color{mundiblue} UNIVERSITÉ MUNDIAPOLIS}\\[0.6cm]
}

{\color{mundigray}\large Business School}\\[0.3cm]
{\color{mundigray}\rule{0.5\textwidth}{0.4pt}}\\[1.0cm]

{\Large\color{mundidark} Cahier d'études de cas}\\[0.35cm]
{\large\color{mundigray} Syllabus M114 -- \S4 Approches pédagogiques}\\[0.8cm]

{\Huge\bfseries\color{mundiblue} Mathématiques appliquées\\[0.15cm] à la gestion}\\[0.9cm]

{\large\itshape Cinq études de cas : seuil de rentabilité,\\
optimisation, capitalisation, amortissement, prévisions}\\[1.6cm]

\begin{tabular}{@{}r l@{}}
\textbf{Filière :} & Licence MGE -- S1 \\[4pt]
\textbf{Module :} & M114 -- Mathématiques appliquées à la gestion \\[4pt]
\textbf{Enseignant(e) :} & Pr Rajae Malek \\[4pt]
\textbf{Année universitaire :} & 2026 -- 2027 \\
\end{tabular}

\vfill
{\color{mundigray}\rule{0.7\textwidth}{0.4pt}}\\[0.3cm]
{\small\color{mundigray} Université Mundiapolis -- Tous droits réservés}
\end{titlepage}
\hypersetup{pageanchor=true}

\tableofcontents
\thispagestyle{fancy}
\clearpage

% =======================================================
\section{Ce que demande le \S4 du syllabus -- et quoi faire}

Le \S4 du syllabus (\emph{Approches pédagogiques utilisées}) n'est \textbf{pas un chapitre à enseigner}, ni une 13\up{e} séance. C'est une \textbf{consigne de méthode} : dans les 12 séances de 3~h, on n'enchaîne pas seulement des exercices techniques. On fait entrer, à des moments précis, une \textbf{étude de cas d'entreprise} que l'étudiant résout comme un gestionnaire, avec des mathématiques.

Les cinq thèmes cités par le syllabus sont :

\begin{center}
\begin{tabular}{@{}p{0.34\textwidth}p{0.28\textwidth}p{0.30\textwidth}@{}}
\toprule
\textbf{Étude de cas} & \textbf{Séance officielle} & \textbf{Outil mathématique} \\
\midrule
Seuil de rentabilité & Séance 5 (continuité / TVI) ; calcul dès le 1\up{er} degré & Équation $\Pi(x)=0$, TVI \\
Optimisation de coût et de profit & Séance 8 (Taylor, extremum) & $f'=0$, $f''$, coût moyen \\
Capitalisation & Séance 3 (suites et finance) & Suite géométrique $(1+i)^n$ \\
Amortissement & Séance 3 (suites et finance) & Suite arithmétique ; somme géométrique \\
Prévisions & Séance 1 (suites de CA) et séance 10 (DL) & Suites ; approximation locale \\
\bottomrule
\end{tabular}
\end{center}

\begin{profbox}[Règle d'or pour les 5 cas]
\begin{enumerate}[leftmargin=1.3cm]
  \item \textbf{Ne pas les regrouper dans une séance unique.} Chacun est branché sur le chapitre où l'outil vient d'être vu.
  \item \textbf{Toujours la même séquence en classe :} (1)~lire l'énoncé d'entreprise $\to$ (2)~traduire en modèle $\to$ (3)~résoudre $\to$ (4)~interpréter en décision (produire / placer / emprunter / prévoir).
  \item \textbf{Interdiction de s'arrêter au nombre.} La dernière phrase de chaque copie doit être une recommandation de gestion.
  \item \textbf{Durée type :} 25 à 45~min par cas, en fin de séance, après le cours. Le corrigé est projeté \emph{après} un temps de recherche, pas avant.
  \item \textbf{Livrable étudiant :} une copie de 1 à 2 pages (calculs + 5 lignes de décision). C'est entraînement à la part « exercices de gestion » de l'examen final (60~\%).
\end{enumerate}
\end{profbox}

\begin{attentionbox}[Ce qu'il ne faut \emph{pas} faire]
Ne pas transformer le cas en cours magistral supplémentaire. L'étudiante ou l'étudiant doit \textbf{écrire le modèle} (profit, suite, dérivée) à partir d'un texte d'entreprise, comme à l'examen. Le professeur circule, débloque le vocabulaire ($CF$, $cv$, annuité, taux), mais ne dicte pas la première équation.
\end{attentionbox}

Le reste de ce cahier donne, pour chaque cas : la fiche professeur (déroulé minute par minute), l'énoncé photocopiable, la méthode, et le corrigé commenté.

% =======================================================
\section{Cas 1 -- Seuil de rentabilité}

\subsection{Fiche professeur}

\begin{profbox}
\textbf{Quand.} Fin de la \textbf{séance 5} (continuité, TVI). La partie A (calcul linéaire) peut aussi servir d'application dès qu'une équation du 1\up{er} degré est disponible.\\
\textbf{Durée.} 40~min (15~min partie A + 15~min partie B + 10~min mise en commun).\\
\textbf{Objectif.} Traduire « à partir de quelle quantité l'entreprise gagne-t-elle de l'argent ? » en $\Pi(x)=0$, puis justifier l'\emph{existence} d'un seuil par le théorème des valeurs intermédiaires, même si la formule n'est pas simple.\\
\textbf{Matériel.} Énoncé ci-dessous (1 feuille). Calculatrice. Tableau pour le graphique $CA$ / $C$.\\
\textbf{Déroulement.}
\begin{enumerate}[leftmargin=1.3cm]
  \item 3~min -- lecture silencieuse. Consigne : « soulignez les données chiffrées et le verbe de décision (produire, arrêter, continuer) ».
  \item 12~min -- partie A en binôme : écrire $C(x)$, $CA(x)$, $\Pi(x)$, calculer $x^\star$ et l'intervalle de profit. Le professeur vérifie que personne n'oublie $x\geqslant 0$ et la capacité.
  \item 12~min -- partie B individuelle : on \emph{interdit} de commencer par chercher les racines. Consigne : « continuité, puis $f(0)$ et $f(4\,000)$, puis conclure \emph{seulement si} les signes sont opposés ». \textbf{Piège prévu :} $f$ reste négative -- le TVI ne s'applique pas, le projet n'est pas rentable sur la capacité. Ne pas le souffler.
  \item 8~min -- mise en commun : un binôme au tableau pour A, un étudiant pour B. Insister : le TVI \textbf{exige un changement de signe} ; sans lui, on ne conclut pas à l'existence d'un seuil.
  \item 5~min -- phrase de décision collective, notée au tableau.
\end{enumerate}
\textbf{Livrable.} Copie individuelle : formules, $x^\star$, intervalle, argument TVI, 5 lignes de recommandation.\\
\textbf{Barème indicatif / 10.} Modèle 2 -- calcul $x^\star$ 2 -- intervalle + capacité 2 -- TVI (hypothèses + conclusion) 3 -- décision 1.
\end{profbox}

\subsection{Énoncé (à photocopier)}

\begin{exobox}[Cas AtlasPack -- seuil de rentabilité]
La PME \textbf{AtlasPack} fabrique des cartons d'emballage. Données du mois :
\begin{itemize}[leftmargin=1.2cm]
  \item coûts fixes $CF = 36\,000$ dh (loyer, encadrement, assurances) ;
  \item coût variable unitaire $cv = 15$ dh ;
  \item prix de vente unitaire $p = 27$ dh ;
  \item capacité maximale de la chaîne : $4\,000$ cartons.
\end{itemize}

\textbf{Partie A -- Calcul du point mort (modèle linéaire)}
\begin{enumerate}[label=\alph*)]
  \item Écrire le coût total $C(x)$, le chiffre d'affaires $CA(x)$ et le profit $\Pi(x)$.
  \item Calculer le seuil de rentabilité $x^\star$ (en quantité) et le seuil en chiffre d'affaires $CA^\star$.
  \item Pour quelles quantités l'entreprise est-elle bénéficiaire, compte tenu de la capacité ?
  \item Les ventes prévues sont $x_0 = 3\,500$ cartons. Calculer la marge de sécurité (en unités et en dh de CA) et conclure.
\end{enumerate}

\textbf{Partie B -- Existence sans formules (théorème des valeurs intermédiaires)}
Le contrôle de gestion fournit une autre structure de coût, plus réaliste (heures supplémentaires, saturation) :
$$C(x) = \frac{x^2}{200} + 8x + 20\,000, \qquad CA(x) = 27x,$$
pour $x\in[0,\,4\,000]$. On pose $f(x)=CA(x)-C(x)$.
\begin{enumerate}[label=\alph*), resume]
  \item Montrer que $f$ est continue sur $[0,\,4\,000]$.
  \item Calculer $f(0)$ et $f(4\,000)$. En déduire, \textbf{sans résoudre} $f(x)=0$, qu'il existe au moins un seuil de rentabilité dans $]0,\,4\,000[$.
  \item Que changerait-on à l'argument si $f(4\,000)$ était encore négatif ?
\end{enumerate}
\end{exobox}

\subsection{Méthode et corrigé}

\begin{methodebox}[Seuil de rentabilité linéaire]
Marge sur coût variable unitaire $m=p-cv$. Tant que $m>0$ :
$$x^\star = \frac{CF}{m} = \frac{CF}{p-cv}, \qquad CA^\star = p\,x^\star = \frac{CF}{m/p}.$$
L'entreprise est bénéficiaire pour $x>x^\star$, à l'intérieur du domaine économique $[0,x_{\max}]$.\\
Marge de sécurité : $x_0-x^\star$ (et $p(x_0-x^\star)$ en CA).
\end{methodebox}

\begin{corrigebox}[Partie A]
\textbf{a)} $C(x)=36\,000+15x$, $CA(x)=27x$, $\Pi(x)=12x-36\,000$.\\
\textbf{b)} $12x^\star-36\,000=0 \Rightarrow x^\star = 3\,000$ cartons. $CA^\star = 27\times 3\,000 = 81\,000$ dh.\\
\textbf{c)} $\Pi(x)>0 \Leftrightarrow x>3\,000$, et $x\leqslant 4\,000$, donc $x\in\,]3\,000,\,4\,000]$. À $x=3\,000$, équilibre (profit nul).\\
\textbf{d)} Marge de sécurité : $3\,500-3\,000=500$ cartons, soit $27\times 500=13\,500$ dh de CA. \textbf{Décision :} le plan à $3\,500$ cartons est rentable, avec une marge confortable ; en dessous de $3\,000$, il faut baisser les fixes, augmenter $p$ ou $m$, ou arrêter.
\end{corrigebox}

\begin{center}
\begin{tikzpicture}[x=0.0017cm, y=0.00004cm]
  \draw[-{Stealth}] (0,0) -- (4600,0) node[right] {\small $x$};
  \draw[-{Stealth}] (0,0) -- (0,120000) node[above] {\small dh};
  \draw[mundiblue, very thick] (0,0) -- (4000,108000) node[pos=0.72, above, sloped] {\small $CA=27x$};
  \draw[mundired, very thick] (0,36000) -- (4000,96000) node[pos=0.82, below, sloped] {\small $C$};
  \draw[dashed, mundigray] (3000,0) -- (3000,81000) -- (0,81000);
  \fill[mundidark] (3000,81000) circle (2.5pt);
  \node[below] at (3000,0) {\small $3\,000$};
  \node[left] at (0,81000) {\small $81\,000$};
\end{tikzpicture}\\
{\small\color{mundigray} Point mort linéaire : intersection de $CA$ et de $C$.}
\end{center}

\begin{corrigebox}[Partie B -- TVI]
\textbf{e)} $f(x)=27x-\bigl(x^2/200+8x+20\,000\bigr)=-x^2/200+19x-20\,000$ est polynomiale, donc continue sur $[0,4\,000]$.\\
\textbf{f)} $f(0)=-20\,000<0$.
$$f(4\,000)=-\frac{16\,000\,000}{200}+76\,000-20\,000=-80\,000+76\,000-20\,000=-24\,000<0.$$
Ici $f(0)<0$ \emph{et} $f(4\,000)<0$ : le TVI \textbf{ne permet pas} de conclure à l'existence d'une racine (les deux bords sont du même signe). On calcule le sommet : $f'(x)=-x/100+19=0 \Rightarrow x=1\,900$, $f(1\,900)=-\,1\,900^2/200+19\times 1\,900-20\,000= -18\,050+36\,100-20\,000=-1\,950<0$. Le maximum est encore négatif : \textbf{aucun} seuil sur $[0,4\,000]$.\\
\textbf{g)} Si $f(4\,000)>0$ (avec $f(0)<0$), le TVI garantirait au moins un $c\in\,]0,4\,000[$ tel que $f(c)=0$. Si $f(4\,000)<0$, deux conclusions possibles : soit le projet n'est jamais rentable sur la capacité (comme ici), soit les racines sont hors intervalle -- d'où l'intérêt de regarder aussi le maximum.
\end{corrigebox}

\begin{attentionbox}[Pourquoi cette partie B est volontairement « piège »]
Beaucoup d'étudiants appliquent le TVI dès qu'on parle de continuité, sans vérifier le \textbf{changement de signe}. Ici les données ont été choisies pour que $f$ reste négative : le cas enseigne à \emph{ne pas} conclure à l'existence d'un seuil par magie. C'est exactement l'esprit de la séance 5.
\end{attentionbox}

\begin{gestionbox}[Variante « TVI qui marche » -- si l'on veut un cas positif]
Mêmes $CA(x)=27x$, mais $C(x)=x^2/200+8x+8\,000$. Alors $f(0)=-8\,000<0$ et $f(4\,000)=-80\,000+76\,000-8\,000=-12\,000$ encore négatif\ldots On prend $C(x)=x^2/400+8x+8\,000$ :
$f(x)=-x^2/400+19x-8\,000$, $f(0)=-8\,000<0$, $f(4\,000)=-40\,000+76\,000-8\,000=28\,000>0$. Le TVI s'applique : il existe un seuil dans $]0,4\,000[$. On peut donner cette variante en devoir.
\end{gestionbox}

% =======================================================
\section{Cas 2 -- Optimisation de coût et de profit}

\subsection{Fiche professeur}

\begin{profbox}
\textbf{Quand.} \textbf{Séance 8} (extremum, conditions du 1\up{er} et du 2\up{nd} ordre), 40 à 45~min en fin de séance.\\
\textbf{Objectif.} Maximiser un profit et minimiser un coût moyen avec $f'=0$ et $f''$, puis traduire « quantité optimale » en décision de production. Retrouver le théorème de gestion : au minimum du coût moyen, \textbf{coût marginal = coût moyen}.\\
\textbf{Déroulement.}
\begin{enumerate}[leftmargin=1.3cm]
  \item 2~min -- rappeler au tableau : $C'(x)=$ coût marginal, $CM(x)=C(x)/x$.
  \item 15~min -- partie A (min $CM$) en binôme. Exiger le tableau de variation ou $CM''$.
  \item 15~min -- partie B (max $\Pi$) individuelle. Contrôler le signe de $\Pi''$.
  \item 8~min -- deux étudiants au tableau + phrase de décision : « produire 10 vs produire 30 n'a pas le même objectif (coût unitaire vs profit total) ».
\end{enumerate}
\textbf{Livrable.} Dérivées, points critiques, justification max/min, valeurs, recommandation.\\
\textbf{Barème / 10.} $CM$ 3 -- lien $C'=CM$ 2 -- $\Pi$ 3 -- $\Pi''$ et décision 2.
\end{profbox}

\subsection{Énoncé (à photocopier)}

\begin{exobox}[Cas Nour Cosmetics -- coût moyen et profit]
\textbf{Nour Cosmetics} produit $x$ lots d'une crème ($x>0$). Le coût total, en dh, est
$$C(x) = x^2 + 20x + 100.$$
Le prix de vente d'un lot est $p=80$ dh (prix imposé par le marché).

\textbf{Partie A -- Minimiser le coût moyen}
\begin{enumerate}[label=\alph*)]
  \item Écrire le coût moyen $CM(x)$ et le coût marginal $C'(x)$.
  \item Déterminer la quantité $x_c$ qui minimise $CM$. Justifier qu'il s'agit d'un minimum.
  \item Vérifier que $C'(x_c)=CM(x_c)$ et interpréter.
\end{enumerate}

\textbf{Partie B -- Maximiser le profit}
\begin{enumerate}[label=\alph*), resume]
  \item Exprimer le profit $\Pi(x)=80x-C(x)$.
  \item Déterminer la quantité $x_\Pi$ qui maximise $\Pi$. Calculer $\Pi_{\max}$.
  \item Comparer $x_c$ et $x_\Pi$. Quelle quantité recommandez-vous, et selon quel critère ?
\end{enumerate}
\end{exobox}

\subsection{Corrigé}

\begin{methodebox}[Optimiser une fonction de gestion]
\begin{enumerate}[leftmargin=1.3cm]
  \item Écrire $f(x)$ (coût moyen ou profit) sur le domaine économique $x>0$.
  \item Calculer $f'(x)=0$ : points critiques.
  \item Signe de $f''$ (ou tableau de variation) : max si $f''<0$, min si $f''>0$.
  \item Revenir aux unités (lots, dh) et formuler une décision.
\end{enumerate}
Au minimum du coût moyen, on a toujours $CM'(x)=0 \Leftrightarrow C'(x)=CM(x)$ (le marginal croise le moyen).
\end{methodebox}

\begin{corrigebox}[Nour Cosmetics]
\textbf{a)} $\displaystyle CM(x)=\frac{x^2+20x+100}{x}=x+20+\frac{100}{x}$, $C'(x)=2x+20$.\\
\textbf{b)} $CM'(x)=1-\dfrac{100}{x^2}=0 \Leftrightarrow x^2=100 \Leftrightarrow x=10$ (car $x>0$).
$CM''(x)=\dfrac{200}{x^3}>0$ pour $x>0$ : minimum. $CM(10)=10+20+10=40$ dh/lot.\\
\textbf{c)} $C'(10)=40=CM(10)$. \textbf{Lecture :} on a intérêt à augmenter la production tant que le coût de la dernière unité est inférieur au coût moyen ; l'optimum unitaire est atteint à l'égalité.\\
\textbf{d)} $\Pi(x)=80x-(x^2+20x+100)=-x^2+60x-100$.\\
\textbf{e)} $\Pi'(x)=-2x+60=0 \Leftrightarrow x=30$. $\Pi''(x)=-2<0$ : maximum.
$\Pi(30)=-900+1\,800-100=800$ dh.\\
\textbf{f)} $x_c=10$ minimise le \emph{coût par lot} ; $x_\Pi=30$ maximise le \emph{profit total}. \textbf{Décision :} si l'objectif est le résultat de l'entreprise, produire $30$ lots (profit $800$ dh). Le palier $x=10$ n'est pertinent que si l'on pilote un coût unitaire (appel d'offres, prix de cession interne), pas le profit.
\end{corrigebox}

% =======================================================
\section{Cas 3 -- Capitalisation}

\subsection{Fiche professeur}

\begin{profbox}
\textbf{Quand.} \textbf{Séance 3} (suites géométriques et finance), 35~min.\\
\textbf{Objectif.} Distinguer intérêt simple / composé ; écrire $C_n=C_0(1+i)^n$ ; calculer un capital futur ; inverser pour trouver une durée (inéquation avec $\ln$ ou essai + encadrement).\\
\textbf{Déroulement.}
\begin{enumerate}[leftmargin=1.3cm]
  \item 5~min -- au tableau, deux colonnes : « simple = suite arithmétique des intérêts » / « composé = suite géométrique du capital ».
  \item 15~min -- calculs a--c en binôme (calculatrice : touche puissance).
  \item 10~min -- question d (durée pour doubler) : d'abord encadrement à la calculatrice $(1{,}05)^n$, puis formule $n>\ln 2/\ln(1{,}i)$.
  \item 5~min -- décision : « pour un placement long, le composé l'emporte ; citer l'écart en dh ».
\end{enumerate}
\textbf{Livrable.} Les quatre résultats numériques + une phrase comparant les deux régimes.\\
\textbf{Barème / 10.} Simple 2 -- composé 3 -- écart 1 -- durée (inéquation + entier) 3 -- décision 1.
\end{profbox}

\subsection{Énoncé (à photocopier)}

\begin{exobox}[Cas Banque Al Aman -- capitalisation]
Un client place $C_0 = 80\,000$ dh à la \textbf{Banque Al Aman}, au taux annuel $i=5\%$, pour une durée de $n$ années, intérêts versés une fois par an.
\begin{enumerate}[label=\alph*)]
  \item \textbf{Intérêt simple.} Écrire le capital $S_n$ au bout de $n$ années. Calculer $S_8$.
  \item \textbf{Intérêt composé.} Justifier que le capital $C_n$ est une suite géométrique ; donner $C_n$ et calculer $C_8$ (arrondir au dh).
  \item Quel régime est plus favorable au client ? Quantifier l'écart au bout de 8 ans.
  \item Combien d'années \emph{entières} faut-il, en intérêt composé, pour que le capital \textbf{double} ? (encadrer puis conclure)
\end{enumerate}
\end{exobox}

\subsection{Corrigé}

\begin{methodebox}[Capitalisation]
Intérêt simple : $S_n=C_0(1+ni)$ (les intérêts ne portent pas intérêt).\\
Intérêt composé : $C_{n}=C_{n-1}(1+i)$, donc $C_n=C_0(1+i)^n$ (suite géométrique de raison $1+i$).\\
Doubler : $(1+i)^n\geqslant 2 \Leftrightarrow n\geqslant \dfrac{\ln 2}{\ln(1+i)}$, puis prendre l'entier supérieur.
\end{methodebox}

\begin{corrigebox}[Banque Al Aman]
\textbf{a)} $S_n=80\,000(1+0{,}05n)$. $S_8=80\,000\times 1{,}4=112\,000$ dh.\\
\textbf{b)} $C_n=80\,000\times(1{,}05)^n$. On a $(1{,}05)^8 \approx 1{,}477455$, d'où $C_8 \approx 118\,196$ dh.\\
\textbf{c)} Le composé est plus favorable : écart $C_8-S_8\approx 6\,196$ dh (les intérêts ont eux-mêmes rapporté des intérêts).\\
\textbf{d)} $(1{,}05)^n\geqslant 2 \Leftrightarrow n\ln(1{,}05)\geqslant \ln 2 \Leftrightarrow n\geqslant \dfrac{\ln 2}{\ln 1{,}05}\approx 14{,}21$.
Contrôle : $(1{,}05)^{14}\approx 1{,}980$, $(1{,}05)^{15}\approx 2{,}079$. Il faut \textbf{15 ans} entiers. (Règle de 72 : $72/5\approx 14{,}4$ ans, cohérente.)\\
\textbf{Décision :} pour 8 ans, le client gagne environ $6\,200$ dh à exiger le composé. Pour un objectif « doubler le capital », communiquer 15 ans, pas 14.
\end{corrigebox}

% =======================================================
\section{Cas 4 -- Amortissement}

\subsection{Fiche professeur}

\begin{profbox}
\textbf{Quand.} Même \textbf{séance 3}, après la capitalisation (les deux cas forment le « bloc finance » du syllabus : \emph{placement, intérêts composés, remboursement d'emprunt}). 45~min au total : 15~min linéaire + 30~min emprunt.\\
\textbf{Objectif.}
\begin{itemize}[leftmargin=1.2cm]
  \item Voir l'amortissement \emph{comptable} d'un équipement comme une \textbf{suite arithmétique} (valeur nette).
  \item Voir le remboursement d'emprunt à annuités constantes comme une \textbf{somme de suite géométrique}.
\end{itemize}
\textbf{Déroulement.}
\begin{enumerate}[leftmargin=1.3cm]
  \item 15~min -- partie A (camion) : très guidée, presque collective. Faire écrire $V_n=V_0-na$.
  \item 5~min -- au tableau, dériver la formule de l'annuité constante :
  $$K=A\cdot\frac{1-(1+i)^{-n}}{i} \;\Longrightarrow\; A=K\cdot\frac{i}{1-(1+i)^{-n}}.$$
  Dire : « chaque annuité rembourse un peu de capital + les intérêts sur le restant dû ».
  \item 20~min -- partie B : calcul de $A$, puis \textbf{deux premières lignes} du tableau d'amortissement (pas les 4 à la main si le temps manque : le reste est en travail personnel).
  \item 5~min -- décision : « l'annuité est-elle soutenable vu un CA annuel de \ldots ? » (donnée dans l'énoncé).
\end{enumerate}
\textbf{Livrable.} $V_n$, $A$ arrondie au dh, tableau (années 1--2 minimum), conclusion de soutenabilité.\\
\textbf{Barème / 12.} Suite arithmétique 3 -- formule et $A$ 4 -- tableau 3 -- décision 2.
\end{profbox}

\subsection{Énoncé (à photocopier)}

\begin{exobox}[Cas TransAtlas -- amortissement d'un camion et d'un emprunt]
\textbf{Partie A -- Amortissement linéaire d'un équipement}\\
TransAtlas achète un camion $V_0=180\,000$ dh, valeur résiduelle $18\,000$ dh au bout de $9$ ans, amortissement linéaire.
\begin{enumerate}[label=\alph*)]
  \item Calculer l'annuité d'amortissement comptable $a$.
  \item Exprimer la valeur nette $V_n$ après $n$ années ($0\leqslant n\leqslant 9$). Quelle est la nature de la suite $(V_n)$ ?
  \item Quelle est la valeur nette au bout de 5 ans ?
\end{enumerate}

\textbf{Partie B -- Remboursement d'emprunt à annuités constantes}\\
Pour financer le camion, la banque prête $K=150\,000$ dh au taux annuel $i=5\%$, remboursable en $n=4$ annuités constantes $A$ (fin d'année).
\begin{enumerate}[label=\alph*), resume]
  \item Calculer $A$ (arrondir au dh).
  \item Compléter le tableau d'amortissement pour les années 1 et 2 (capital restant en début, intérêts, amortissement du capital, annuité, capital restant en fin).
  \item Le chiffre d'affaires annuel prévisionnel lié au camion est $90\,000$ dh. L'annuité est-elle soutenable ? Argumenter.
\end{enumerate}
\end{exobox}

\subsection{Corrigé}

\begin{methodebox}[Deux amortissements à ne pas confondre]
\textbf{Comptable (linéaire) :} on répartit $V_0-V_{\text{résiduelle}}$ en $N$ parts égales. Suite \textbf{arithmétique} de raison $-a$.\\
\textbf{Financier (emprunt) :} on verse chaque année la même annuité $A$ qui couvre intérêts + capital. Les intérêts décroissent, la part capital croît. Formule :
$$A = K\,\frac{i}{1-(1+i)^{-n}}.$$
Intérêts de l'année $k$ : $i\times$ (restant dû en début). Capital remboursé : $A-$ intérêts. Le dernier restant dû doit être $\approx 0$ (arrondis).
\end{methodebox}

\begin{corrigebox}[Partie A -- camion]
\textbf{a)} $a=(180\,000-18\,000)/9=18\,000$ dh/an.\\
\textbf{b)} $V_n=180\,000-18\,000\,n$. Suite arithmétique de premier terme $180\,000$ et de raison $-18\,000$.\\
\textbf{c)} $V_5=180\,000-90\,000=90\,000$ dh.
\end{corrigebox}

\begin{corrigebox}[Partie B -- emprunt]
\textbf{d)} $(1{,}05)^4=1{,}21550625$, $(1{,}05)^{-4}\approx 0{,}8227025$,
$$A=150\,000\times\frac{0{,}05}{1-(1{,}05)^{-4}} \approx \frac{7\,500}{0{,}1772975} \approx 42\,302 \text{ dh.}$$
\textbf{e)} Tableau (montants arrondis au dh) :
\begin{center}
\small
\begin{tabular}{crrrrr}
\toprule
Année & Début & Intérêts & Capital & Annuité & Fin \\
\midrule
1 & $150\,000$ & $7\,500$ & $34\,802$ & $42\,302$ & $115\,198$ \\
2 & $115\,198$ & $5\,760$ & $36\,542$ & $42\,302$ & $78\,656$ \\
\bottomrule
\end{tabular}
\end{center}
(Année 1 : $150\,000\times 0{,}05=7\,500$, $42\,302-7\,500=34\,802$, $150\,000-34\,802=115\,198$. Année 2 : $115\,198\times 0{,}05\approx 5\,760$.)\\
\textbf{f)} $A\approx 42\,302$ représente $47\%$ du CA annuel $90\,000$. \textbf{Décision :} l'annuité est lourde : il reste $47\,698$ dh de CA pour les charges d'exploitation du camion (carburant, chauffeur, assurance). On recommande de vérifier la marge après ces charges ; si elle est mince, allonger $n$ (baisse $A$) ou augmenter l'apport pour réduire $K$.
\end{corrigebox}

% =======================================================
\section{Cas 5 -- Prévisions}

\subsection{Fiche professeur}

\begin{profbox}
\textbf{Quand.} En deux temps, comme le planning M114.
\begin{itemize}[leftmargin=1.2cm]
  \item \textbf{Séance 1} (25~min) : prévision par suite de chiffre d'affaires (partie A).
  \item \textbf{Séance 10} (20~min) : prévision locale par approximation (DL / coût marginal, partie B).
\end{itemize}
\textbf{Objectif.} Faire une prévision \emph{chiffrée} et \emph{discuter le modèle} (arithmétique vs géométrique ; approximation vs valeur exacte). Le syllabus parle d'« identifier des tendances » : c'est ici.\\
\textbf{Déroulement séance 1.}
\begin{enumerate}[leftmargin=1.3cm]
  \item 5~min -- « une prévision n'est pas une vérité, c'est un modèle ». Deux hypothèses au tableau.
  \item 12~min -- calculs des CA 2028 sous A et sous B.
  \item 8~min -- débat : quelle hypothèse pour une PME en croissance ? On exige une phrase de prudence (intervalle, pas un seul chiffre).
\end{enumerate}
\textbf{Déroulement séance 10.} 15~min de calcul (exact vs $C'(x_0)h$) + 5~min : « le marginal sert à prévoir un \emph{petit} choc de production ».\\
\textbf{Livrable.} Deux prévisions + choix motivé (A) ; erreur relative de l'approximation (B).\\
\textbf{Barème / 10.} Suite arithmétique 2 -- géométrique 3 -- choix du modèle 2 -- approximation et erreur 3.
\end{profbox}

\subsection{Énoncé (à photocopier)}

\begin{exobox}[Cas Maison du Thé -- prévisions]
Le chiffre d'affaires 2025 de la \textbf{Maison du Thé} est $CA_0=1\,200\,000$ dh. La direction demande une prévision pour \textbf{2028} (soit 3 ans plus tard). Deux hypothèses circulent :
\begin{itemize}[leftmargin=1.2cm]
  \item \textbf{Hypothèse A :} hausse de $100\,000$ dh chaque année (suite arithmétique).
  \item \textbf{Hypothèse B :} hausse de $7\%$ chaque année (suite géométrique).
\end{itemize}

\textbf{Partie A -- Prévoir un CA (séance 1)}
\begin{enumerate}[label=\alph*)]
  \item Écrire $CA_n$ dans chaque hypothèse ($n=0$ en 2025).
  \item Calculer les deux prévisions pour 2028.
  \item Laquelle présentez-vous au conseil d'administration, et sous quelle réserve ?
\end{enumerate}

\textbf{Partie B -- Prévoir un petit choc de production (séance 10)}
Le coût mensuel d'un atelier est $C(x)=x^2+20x+100$ (en dh, $x$ en lots). On produit actuellement $x_0=30$ lots.
\begin{enumerate}[label=\alph*), resume]
  \item Calculer $C(30)$ et $C(32)$ (hausse de 2 lots). En déduire le surcoût exact.
  \item Approcher le surcoût par $C'(30)\times 2$. Calculer l'erreur absolue et l'erreur relative.
  \item Pourquoi cette approximation est-elle un outil de \emph{prévision} pour le contrôle de gestion ?
\end{enumerate}
\end{exobox}

\subsection{Corrigé}

\begin{methodebox}[Prévisions]
Suite arithmétique (croissance en \emph{niveau}) : $u_n=u_0+nd$.\\
Suite géométrique (croissance en \emph{taux}) : $u_n=u_0\,q^n$.\\
Petit choc : $f(x_0+h)\approx f(x_0)+f'(x_0)h$ (DL d'ordre 1 / coût marginal). Toujours comparer à la valeur exacte et donner l'erreur.
\end{methodebox}

\begin{corrigebox}[Partie A]
\textbf{a)} A : $CA_n=1\,200\,000+100\,000\,n$. B : $CA_n=1\,200\,000\times(1{,}07)^n$.\\
\textbf{b)} 2028 $\Leftrightarrow n=3$. A : $1\,500\,000$ dh. B : $1\,200\,000\times(1{,}07)^3=1\,200\,000\times 1{,}225043=1\,470\,052$ dh (arrondi).\\
\textbf{c)} Les deux chiffres sont proches (écart $30\,000$ dh, $2\%$). \textbf{Décision :} présenter un \emph{intervalle} $[1{,}47$ M$;\,1{,}50$ M$]$ plutôt qu'un point. Hypothèse B (taux) si l'on croit à une croissance proportionnelle au CA ; hypothèse A si le marché est saturé et que l'on ajoute un même volume chaque année. Mentionner le risque (aléa climatique, concurrence).
\end{corrigebox}

\begin{corrigebox}[Partie B]
\textbf{d)} $C(30)=900+600+100=1\,600$ dh, $C(32)=1\,024+640+100=1\,764$ dh. Surcoût exact : $164$ dh.\\
\textbf{e)} $C'(x)=2x+20$, $C'(30)=80$, approximation $80\times 2=160$ dh. Erreur absolue $4$ dh, erreur relative $4/164 \approx 2{,}4\%$.\\
\textbf{f)} Le contrôleur n'a souvent sous la main que le coût marginal (ou une estimation). Pour une petite variation de production, $C'(x_0)h$ donne une prévision rapide du surcoût, ici à $2\%$ près. Si $h$ devient grand, l'approximation se dégrade (terme $h^2$ du Taylor) : il faut recalculer $C(x_0+h)$.
\end{corrigebox}

% =======================================================
\section{Calendrier d'insertion dans les 12 séances}

\begin{center}
\renewcommand{\arraystretch}{1.15}
\begin{tabular}{@{}clp{0.28\textwidth}p{0.42\textwidth}@{}}
\toprule
\textbf{Séance} & \textbf{Vol.} & \textbf{Cas à faire} & \textbf{Moment dans les 3~h} \\
\midrule
1 & 3~h & Prévisions, partie A (Maison du Thé) & 25~min en fin, après les suites de CA / stocks \\
3 & 3~h & Capitalisation + Amortissement & 35~min + 45~min après le cours sur les suites particulières \\
5 & 3~h & Seuil de rentabilité (AtlasPack) & 40~min après le TVI \\
8 & 3~h & Optimisation (Nour Cosmetics) & 45~min après extremum / Taylor \\
10 & 3~h & Prévisions, partie B (choc de coût) & 20~min après les équivalents / DL \\
\bottomrule
\end{tabular}
\end{center}

Les séances 2, 4, 6, 7, 9, 11, 12 restent en TD classiques (démonstrations, calculs, annales). Le contrôle continu (séance 6 ou 7) peut reprendre une question courte de type « seuil » ou « suite de CA », pas un cas entier.

\begin{propbox}[Ce que l'examen final doit refléter]
Le syllabus précise que l'écrit final (60~\%) « couvre l'ensemble du programme, avec une part d'exercices de gestion ». Au moins \textbf{un exercice sur 3 ou 4} de l'examen doit être un mini-cas du même format : texte d'entreprise, 3 à 5 questions, dernière question = décision. Les cinq cas de ce cahier sont le vivier d'annales.
\end{propbox}

\section*{Checklist professeur (à imprimer au verso du planning)}
\begin{itemize}[leftmargin=1.2cm]
  \item [\,] Séance 1 : photocopies Maison du Thé (partie A seulement).
  \item [\,] Séance 3 : photocopies Al Aman + TransAtlas ; calculatrices ; formule de $A$ au tableau.
  \item [\,] Séance 5 : photocopies AtlasPack ; prévoir la variante TVI « qui marche » en devoir si la partie B piège a trop déstabilisé.
  \item [\,] Séance 8 : photocopies Nour Cosmetics ; rappeler $CM$ vs $\Pi$ avant de lancer.
  \item [\,] Séance 10 : reprendre la feuille Maison du Thé, partie B.
  \item [\,] Après chaque cas : 5 lignes de décision notées ; une copie ramassée au hasard pour réguler le rythme.
\end{itemize}

\vspace{0.5cm}
{\small\color{mundigray}\textit{Document d'accompagnement du syllabus M114 \S4 -- à utiliser en complément du support de cours, non à la place des chapitres théoriques.}}

\end{document}
